\chapter{{\v C}ech computation of higher direct images $R^i f_*$ (unconditional)}
\label{chap:Cohomology_CechHigherDirectImage}
% archon:covers AlgebraicJacobian/Cohomology/CechHigherDirectImage.lean
% SOURCE: [Stacks Project], Cohomology of Schemes, \S{\v C}ech cohomology of
% quasi-coherent sheaves and \S Quasi-coherence of higher direct images
% (Lemmas \texttt{lemma-cech-cohomology-quasi-coherent-trivial},
% \texttt{lemma-quasi-coherent-affine-cohomology-zero},
% \texttt{lemma-cech-cohomology-quasi-coherent},
% \texttt{lemma-quasi-coherence-higher-direct-images-application},
% \texttt{lemma-flat-base-change-cohomology}) (read from
% references/stacks-coherent.tex).

\section{Setup and motivation}

This chapter constructs the higher derived direct images \(R^i f_*\mathcal{F}\) for
\(i \geq 1\) \emph{without appealing to injective resolutions in the category of
sheaves of modules}. The companion development of
Chapter~\ref{chap:Cohomology_HigherDirectImage} defines \(R^i f_*\) as a right
derived functor, which in a formalisation requires the ambient category of
\(\mathcal{O}_X\)-modules to have enough injectives; that property is not currently
available for sheaves of modules over a sheaf of rings whose value category varies
over the site. The {\v C}ech approach developed here sidesteps the issue entirely:
it computes \(R^i f_*\mathcal{F}\) as the cohomology of an explicit complex built
from the pushforwards of \(\mathcal{F}\) over the finite intersections of an affine
open cover, and so produces an \emph{unconditional} construction of \(R^i f_*\) for
quasi-coherent \(\mathcal{F}\) and separated quasi-compact \(f\).

Throughout, \(f : X \to S\) is a morphism of schemes that is quasi-compact and
separated (so that all finite intersections of an affine open cover of \(X\) are
again affine), and \(\mathcal{F}\) is a quasi-coherent \(\mathcal{O}_X\)-module. We
write \(\operatorname{QCoh}(X)\) for the category of quasi-coherent
\(\mathcal{O}_X\)-modules. A base change of \(f\) along a morphism \(g : S' \to S\)
is recorded by the cartesian square obtained from the fibre product
\(X' = X \times_S S'\):
\[
  \begin{array}{ccc}
    X' & \xrightarrow{\;g'\;} & X \\[2pt]
    {\scriptstyle f'}\big\downarrow & & \big\downarrow{\scriptstyle f} \\[2pt]
    S' & \xrightarrow{\;g\;} & S
  \end{array}
\]
Here \(f' : X' \to S'\) is the base change of \(f\), \(g' : X' \to X\) is the first
projection, and \(\mathcal{F}' = (g')^*\mathcal{F}\) is the pullback of
\(\mathcal{F}\) to \(X'\).

\section{The {\v C}ech nerve and {\v C}ech complex of an affine cover}

\begin{definition}

\leanok
[{\v C}ech nerve of an affine open cover]
  \label{def:cech_nerve}
  \lean{AlgebraicGeometry.CechNerve}
  \uses{def:cover_cech_nerve, lem:push_pull_functor}
  % SOURCE: [Stacks Project], Cohomology of Schemes, \S{\v C}ech cohomology of
  % quasi-coherent sheaves, opening paragraph and Lemma
  % \texttt{lemma-cech-cohomology-quasi-coherent-trivial}
  % (read from references/stacks-coherent.tex, L36--67).
  % SOURCE QUOTE: "Let $X$ be a scheme.
  % Let $U \subset X$ be an affine open.
  % Recall that a {\it standard open covering} of $U$ is a covering
  % of the form $\mathcal{U} : U = \bigcup_{i = 1}^n D(f_i)$
  % where $f_1, \ldots, f_n \in \Gamma(U, \mathcal{O}_X)$ generate
  % the unit ideal, see
  % Schemes, Definition \ref{schemes-definition-standard-covering}."
  \textit{Source: Stacks Project, Cohomology of Schemes, \S{\v C}ech cohomology of
  quasi-coherent sheaves.}
  Let \(X\) be a scheme and let
  \(\mathcal{U} : X = \bigcup_{i \in I} U_i\) be a finite open cover of \(X\) by
  affine opens, with index set \(I\). For a finite nonempty tuple
  \(i_0, \ldots, i_p \in I\) write
  \(U_{i_0 \ldots i_p} = U_{i_0} \cap \cdots \cap U_{i_p}\) for the corresponding
  intersection. The \emph{{\v C}ech nerve} of \(\mathcal{U}\) is the augmented
  \emph{cosimplicial} object in \(\operatorname{QCoh}(X)\) whose object in
  cosimplicial degree \(p\) is the direct image, along the open immersions
  \(j_{i_0 \ldots i_p} : U_{i_0 \ldots i_p} \hookrightarrow X\), of the restriction
  of \(\mathcal{O}_X\) (or of any chosen quasi-coherent sheaf) to the
  \((p+1)\)-fold intersections,
  \[
    [p] \;\longmapsto\;
      \prod_{(i_0, \ldots, i_p) \in I^{p+1}}
        (j_{i_0 \ldots i_p})_* \,\bigl(\,\cdot\,|_{U_{i_0 \ldots i_p}}\bigr),
  \]
  with cofaces given by the restriction maps that omit one index and
  codegeneracies by the maps that repeat one index; the augmentation in degree
  \(-1\) is the object on all of \(X\). The variance is essential: the geometric
  nerve of the cover (the iterated fibre powers of \(\coprod_i U_i\) over \(X\),
  introduced in \S\ref{sec:cech_three_part} below) is a \emph{simplicial} object
  in \emph{schemes}, but applying the covariant direct-image functor turns it into
  a \emph{cosimplicial} object in \(\operatorname{QCoh}(X)\); consequently the
  associated alternating-coface complex is a \emph{cochain} complex (differentials
  raising degree), not a chain complex. When \(X\) is separated each intersection
  \(U_{i_0 \ldots i_p}\) is itself affine. The local affine model is the
  \emph{standard open covering} \(U = \bigcup_{i=1}^n D(f_i)\) of an affine open
  \(U = \operatorname{Spec}(A)\) by basic opens, where \(f_1, \ldots, f_n \in A\)
  generate the unit ideal; on such a cover the intersections are again basic opens
  \(D(f_{i_0} \cdots f_{i_p})\).
\end{definition}

\begin{definition}

\leanok
[{\v C}ech complex of a quasi-coherent sheaf]
  \label{def:cech_complex}
  \lean{AlgebraicGeometry.CechComplex}
  \uses{def:cech_nerve, def:relative_cech_complex_of_nerve}
  % SOURCE: [Stacks Project], Cohomology of Schemes, Lemma
  % \texttt{lemma-cech-cohomology-quasi-coherent-trivial} (proof)
  % (read from references/stacks-coherent.tex, L59--77).
  % SOURCE QUOTE: "Clearly the {\v C}ech complex
  % $\check{\mathcal{C}}^\bullet(\mathcal{U}, \mathcal{F})$
  % is identified with the complex
  % $$
  % \prod\nolimits_{i_0} M_{f_{i_0}} \to
  % \prod\nolimits_{i_0i_1} M_{f_{i_0}f_{i_1}} \to
  % \prod\nolimits_{i_0i_1i_2} M_{f_{i_0}f_{i_1}f_{i_2}} \to
  % \ldots
  % $$
  % We are asked to show that the extended complex
  % \begin{equation}
  % \label{equation-extended}
  % 0 \to
  % M \to
  % \prod\nolimits_{i_0} M_{f_{i_0}} \to
  % \prod\nolimits_{i_0i_1} M_{f_{i_0}f_{i_1}} \to
  % \prod\nolimits_{i_0i_1i_2} M_{f_{i_0}f_{i_1}f_{i_2}} \to
  % \ldots
  % \end{equation}
  % (whose truncation we have studied in
  % Algebra, Lemma \ref{algebra-lemma-cover-module}) is exact."
  \textit{Source: Stacks Project, Cohomology of Schemes,
  \texttt{lemma-cech-cohomology-quasi-coherent-trivial} (standard-cover {\v C}ech
  vanishing).}
  Let \(\mathcal{F}\) be a quasi-coherent \(\mathcal{O}_X\)-module and let
  \(\mathcal{U} : X = \bigcup_{i \in I} U_i\) be a finite affine open cover with
  all intersections \(U_{i_0 \ldots i_p}\) affine (e.g.\ \(X\) separated). The
  \emph{{\v C}ech complex} \(\check{\mathcal{C}}^\bullet(\mathcal{U}, \mathcal{F})\)
  is the cochain complex obtained by applying the global-sections functor to the
  {\v C}ech nerve of Definition~\ref{def:cech_nerve} evaluated at \(\mathcal{F}\);
  in degree \(p\) it is
  \[
    \check{\mathcal{C}}^p(\mathcal{U}, \mathcal{F})
      = \prod_{(i_0, \ldots, i_p) \in I^{p+1}}
        \mathcal{F}\bigl(U_{i_0 \ldots i_p}\bigr),
  \]
  with differential \(d : \check{\mathcal{C}}^p \to \check{\mathcal{C}}^{p+1}\) the
  alternating sum \((d s)_{i_0 \ldots i_{p+1}} = \sum_{j=0}^{p+1} (-1)^j\,
  s_{i_0 \ldots \widehat{i_j} \ldots i_{p+1}}|_{U_{i_0 \ldots i_{p+1}}}\) of the
  restriction maps. Over an affine \(U = \operatorname{Spec}(A)\) with
  \(\mathcal{F}|_U = \widetilde{M}\) and a standard cover by the
  \(D(f_i)\), this is the complex
  \(\prod_{i_0} M_{f_{i_0}} \to \prod_{i_0 i_1} M_{f_{i_0} f_{i_1}} \to \cdots\) of
  localisations. The \emph{relative {\v C}ech complex} \(\check{\mathcal{C}}^\bullet
  (\mathfrak{U}, \mathcal{F})\) for \(f : X \to S\) is the analogous complex in
  \(\operatorname{QCoh}(S)\) obtained by applying \(f_*\) over each finite
  intersection instead of global sections; its terms are
  \(\prod_{i_0 \ldots i_p} (f|_{U_{i_0 \ldots i_p}})_*
  \bigl(\mathcal{F}|_{U_{i_0 \ldots i_p}}\bigr)\), and it is a complex in
  \(\operatorname{QCoh}(S)\) because each \(U_{i_0 \ldots i_p}\) is affine and the
  pushforward of a quasi-coherent sheaf along a quasi-compact quasi-separated
  morphism is quasi-coherent.
\end{definition}

\subsection{The three-part construction of the relative {\v C}ech complex}
\label{sec:cech_three_part}

It is convenient to factor the construction of
Definitions~\ref{def:cech_nerve}--\ref{def:cech_complex} into three independent
pieces. Two of them are purely formal and unconditional; the entire mathematical
content that is not yet available off the shelf is isolated in the middle piece.

\paragraph{(1) Geometric backbone — the augmented {\v C}ech nerve of an arrow.}
Package the affine cover as a single morphism
\(\coprod_{i \in I} U_i \longrightarrow X\) (the canonical map out of the
coproduct induced by the inclusions \(U_i \hookrightarrow X\)). The augmented
{\v C}ech nerve of this arrow is an augmented \emph{simplicial scheme} over
\(X\): its object in simplicial degree \(p\) is the \((p+1)\)-fold fibre power
\[
  \Bigl(\coprod_{i} U_i\Bigr)^{\times_X (p+1)}
    \;=\;
    \coprod_{(i_0, \ldots, i_p) \in I^{p+1}}
      U_{i_0} \times_X \cdots \times_X U_{i_p},
\]
with the augmentation the cover map to \(X\). Because the \(U_i\) are open
subschemes of \(X\), each fibre product \(U_{i_0} \times_X \cdots \times_X U_{i_p}\)
is canonically the intersection \(U_{i_0 \ldots i_p}\). This backbone exists
\emph{unconditionally}: it uses only that the category of schemes has coproducts
and finite limits (hence the wide pullbacks underlying the {\v C}ech nerve of an
arrow). It carries no sheaf-theoretic data yet — it is the simplicial diagram of
spaces over which the construction is indexed.

\paragraph{(2) Push--pull functor — lifting the backbone to modules.}
The geometric backbone is turned into a cosimplicial \(\mathcal{O}_X\)-module by
the \emph{relative-direct-image functor on the over-category}
\[
  G : (X/\mathbf{Sch})^{\mathrm{op}} \longrightarrow \operatorname{QCoh}(X),
  \qquad
  (Y \xrightarrow{\,p\,} X) \;\longmapsto\; p_*\, p^*\mathcal{F},
\]
sending an \(X\)-scheme \(p : Y \to X\) to the pushforward along \(p\) of the
pullback \(p^*\mathcal{F}\). Composing \(G\) with the (opposite of the) simplicial
backbone of (1) produces exactly the augmented cosimplicial object of
Definition~\ref{def:cech_nerve}: on the degree-\(p\) piece
\(p : U_{i_0 \ldots i_p} \hookrightarrow X\) one has
\(p_* p^*\mathcal{F} = (j_{i_0 \ldots i_p})_*(\mathcal{F}|_{U_{i_0 \ldots i_p}})\),
and a morphism of \(X\)-schemes \(h : (Y_1, p_1) \to (Y_2, p_2)\) (so
\(h \circ p_2 = p_1\)\,, reading \(p\) as the structure map) induces the
restriction comparison \(G(p_2) \to G(p_1)\) via the unit/counit of the
pull-push adjunction together with the over-triangle identifying the two
structure maps. This functor \(G\) is the \emph{one non-formal step} in the whole
construction. Its functoriality --- that \(G\) respects identities and composition
of morphisms over \(X\) --- is precisely a statement about the coherence of the
comparison isomorphisms \((p \circ q)_* \cong p_* q_*\) and
\((p \circ q)^* \cong q^* p^*\) for pushforward and pullback of sheaves of
modules, together with their compatibility with the adjunction units. This
functoriality is \emph{not} coupled to the project's tensor--pullback substrate;
it consumes only Mathlib's pseudofunctor packaging of pullback and pushforward.
Concretely, pullback and pushforward of sheaves of modules are organised as a
pseudofunctor \(\mathrm{LocallyDiscrete}(\mathbf{Sch}^{\mathrm{op}}) \to
\mathbf{Adj}\,\mathbf{Cat}\) (sending a morphism to the adjoint pair
\(p^* \dashv p_*\)), and the comparison isomorphisms
\((p \circ q)^* \cong q^* p^*\), \((p \circ q)_* \cong p_* q_*\) are its structure
\(2\)-cells. The functoriality of \(G\) needs exactly three facts about this
packaging: the mate identity expressing the conjugate of the inverse pullback
comparison as the pushforward comparison (the lemma
\texttt{conjugateEquiv\_pullbackComp\_inv}), its identity normalisation
\texttt{conjugateEquiv\_pullbackId\_hom}, and the pseudofunctor unitor and pentagon
(associativity) coherences \texttt{pseudofunctor\_left\_unitality},
\texttt{pseudofunctor\_right\_unitality}, \texttt{pseudofunctor\_associativity}.
All of these are already available off the shelf, so the functoriality of the
relative-direct-image functor is a \emph{consumer of Mathlib's coherence},
independent of the project-local sheafification/pullback machinery of the
tensor--pullback substrate. The laws are nonetheless not formal trivialities: each
is a multi-step mate calculation that transports the comparison isomorphisms along
the over-triangle \(g \circ p_2 = p_1\) of structure maps and past the adjunction
unit, gluing them across two \(\mathrm{eqToHom}\) coercions; the obstacle is the
bookkeeping of those transports, not a missing geometric ingredient.

\paragraph{(3) Coherence-free plumbing — from the nerve to the cochain complex.}
The passage from the augmented cosimplicial \(\mathcal{O}_X\)-module of (1)--(2)
to the relative {\v C}ech complex in \(\operatorname{QCoh}(S)\) is again purely
formal and uses \emph{no} pushforward/pullback coherence. One forgets the
augmentation, pushes the cosimplicial object forward along \(f : X \to S\) (a
covariant operation applied degreewise), and takes the alternating-coface-map
cochain complex of the resulting cosimplicial object of \(\mathcal{O}_S\)-modules.
This step relies only on the preadditivity of \(\operatorname{QCoh}(S)\): the
alternating sum \(\sum_{j} (-1)^j d^j\) of the cofaces is defined in any
preadditive category and squares to zero by the cosimplicial identities. Thus the
relative {\v C}ech complex \(\check{\mathcal{C}}^\bullet(\mathfrak{U}, \mathcal{F})\)
of Definition~\ref{def:cech_complex} is genuinely \emph{defined} in terms of the
{\v C}ech nerve of Definition~\ref{def:cech_nerve}: once the nerve is constructed,
the complex follows with no further input. In summary, of the three pieces only
the push--pull functor (2) requires content beyond formal category theory, and
that content is the pushforward/pullback pseudofunctor coherence supplied by
Mathlib (the comparison \(2\)-cells of the pseudofunctor
\(\mathrm{LocallyDiscrete}(\mathbf{Sch}^{\mathrm{op}}) \to \mathbf{Adj}\,\mathbf{Cat}\)
and their unitor/pentagon identities), not any project-local ingredient.

\paragraph{Formal bricks of the three-part construction.}
The geometric backbone (1), the push--pull functor (2), and the plumbing (3) are
recorded as the following declarations. The backbone and plumbing are
unconditional; the push--pull object and morphism maps are likewise unconditional
bricks, and the functor laws assembling them are the lone coherence statement.

\begin{definition}\leanok
[Arrow of an affine cover]
  \label{def:cover_arrow}
  \lean{AlgebraicGeometry.coverArrow}
  Let \(\mathcal{U} : X = \bigcup_{i \in I} U_i\) be an affine open cover. Its
  \emph{cover arrow} is the single morphism
  \(\coprod_{i \in I} U_i \longrightarrow X\) obtained as the universal map out of
  the coproduct induced by the open immersions \(U_i \hookrightarrow X\). Packaging
  the cover as one arrow lets the generic {\v C}ech-nerve-of-an-arrow machinery
  apply.
\end{definition}

\begin{definition}\leanok
[Augmented {\v C}ech nerve of a cover]
  \label{def:cover_cech_nerve}
  \lean{AlgebraicGeometry.coverCechNerve}
  \uses{def:cover_arrow}
  The augmented {\v C}ech nerve of the cover arrow of
  Definition~\ref{def:cover_arrow}: the augmented \emph{simplicial scheme} over
  \(X\) whose degree-\(p\) object is the \((p+1)\)-fold fibre power
  \[
    \Bigl(\coprod_{i} U_i\Bigr)^{\times_X (p+1)}
      = \coprod_{(i_0, \ldots, i_p) \in I^{p+1}}
        U_{i_0} \times_X \cdots \times_X U_{i_p},
  \]
  with augmentation the cover map to \(X\). It exists unconditionally because the
  category of schemes has all finite limits (hence the wide pullbacks underlying the
  {\v C}ech nerve of an arrow). This is the geometric backbone (1).
\end{definition}

\begin{definition}\leanok
[Push--pull object map]
  \label{def:push_pull_obj}
  \lean{AlgebraicGeometry.pushPullObj}
  The object map of the push--pull functor \(G\): to an \(X\)-scheme
  \(p : Y \to X\) it assigns
  \[
    G(Y, p) \;=\; p_*\, p^*\mathcal{F},
  \]
  the pushforward along \(p\) of the pullback \(p^*\mathcal{F}\).
\end{definition}

\begin{definition}\leanok
[Push--pull comparison map]
  \label{def:push_pull_map}
  \lean{AlgebraicGeometry.pushPullMap}
  \uses{def:push_pull_obj}
  The morphism map of \(G\). A morphism \(g : (Y_2, p_2) \to (Y_1, p_1)\) of
  \(X\)-schemes has underlying map \(\bar g : Y_2 \to Y_1\) satisfying the
  over-triangle \(p_1 \circ \bar g = p_2\); the contravariant functor produces the
  restriction comparison \(G(Y_1, p_1) \to G(Y_2, p_2)\) as the five-step composite
  \[
    p_{1*}\,p_1^*\mathcal{F}
      \xrightarrow{\,p_{1*}(\eta)\,}
    p_{1*}\,\bar g_*\,\bar g^*\,p_1^*\mathcal{F}
      \xrightarrow{\,\sim\,}
    (p_1 \circ \bar g)_*\,\bar g^*\,p_1^*\mathcal{F}
      \;=\;
    p_{2*}\,\bar g^*\,p_1^*\mathcal{F}
      \xrightarrow{\,p_{2*}(\sim)\,}
    p_{2*}\,p_2^*\mathcal{F},
  \]
  where \(\eta\) is the unit of the adjunction \(\bar g^* \dashv \bar g_*\), the
  second arrow is the pushforward comparison \(p_{1*}\bar g_* \cong (p_1 \bar g)_*\),
  the equality is the substitution \(p_2 = p_1 \circ \bar g\) of the over-triangle,
  and the last arrow applies \(p_{2*}\) to the pullback comparison
  \(\bar g^* p_1^* \cong (p_1 \bar g)^* = p_2^*\). The two over-triangle
  substitutions are realised as coercions along \(p_1 \circ \bar g = p_2\).
\end{definition}

\begin{lemma}\leanok
[Push--pull functor assembly]
  \label{lem:push_pull_functor}
  \lean{AlgebraicGeometry.pushPullMap_id}
  \lean{AlgebraicGeometry.pushPullMap_comp}
  % NOTE (iter-264): only `pushPullMap_id` exists in Lean (landed axiom-clean this iter);
  % `pushPullMap_comp` (the pentagon law) is still deferred to an in-file comment, NOT a
  % declaration. The statement-block \leanok therefore over-states this lemma. Plan agent:
  % either split this block into two (one \lean{} pin each) so \leanok tracks them
  % independently, or have a `lemma pushPullMap_comp ... := sorry` stub added to the Lean file.
  % NOTE iter-283: removed the spurious `def:cech_nerve` from this \uses. The
  % functor laws `pushPullMap_id`/`pushPullMap_comp` (L216/278 of the Lean file)
  % depend only on push_pull_obj/push_pull_map and the mate lemmas; their
  % signatures never mention `CechNerve`. The dependency runs the OTHER way:
  % `CechNerve` is BUILT FROM the push-pull functor (def:cech_nerve \uses
  % lem:push_pull_functor), so listing def:cech_nerve here created a
  % cech_nerve <-> push_pull_functor cycle that crashed `leanblueprint web`.
  \uses{def:push_pull_obj, def:push_pull_map, lem:push_pull_unit_mate, lem:push_pull_transport_cancel}
  The object and morphism assignments of
  Definitions~\ref{def:push_pull_obj}--\ref{def:push_pull_map} assemble into a
  functor \(G : (X/\mathbf{Sch})^{\mathrm{op}} \to \operatorname{QCoh}(X)\): the
  comparison maps respect identities and composition,
  \[
    G(\mathrm{id}_Y) = \mathrm{id}_{G(Y)},
    \qquad
    G(g \circ h) = G(h) \circ G(g),
  \]
  the latter for composable morphisms \(h : Y_3 \to Y_2\), \(g : Y_2 \to Y_1\) of
  \(X\)-schemes. Composing \(G\) with the backbone of
  Definition~\ref{def:cover_cech_nerve} yields the {\v C}ech nerve of
  Definition~\ref{def:cech_nerve}.
\end{lemma}

\begin{proof}
  \uses{def:push_pull_obj, def:push_pull_map}
  Both laws are statements purely about the pushforward/pullback comparison
  isomorphisms and carry no sectionwise content. The route is the adjunction-mate
  calculus applied to the head \(p_{1*}(\eta)\) followed by the pushforward
  comparison in the composite defining \(G(g)\): rewriting this head as the
  conjugate-mate of the inverse pullback comparison
  \((p \circ q)^* \cong q^* p^*\) — via \texttt{conjugateEquiv\_pullbackComp\_inv} —
  turns both laws into identities among the pseudofunctor's structure \(2\)-cells.

  For \(G(\mathrm{id}_Y) = \mathrm{id}\): after the mate rewrite the identity leg
  collapses by the identity normalisation \texttt{conjugateEquiv\_pullbackId\_hom}
  together with the right unitality \texttt{pseudofunctor\_right\_unitality};
  the two over-triangle coercions cancel against the unitality coercion.

  For \(G(g \circ h) = G(h) \circ G(g)\): the same mate rewrite reduces the claim to
  reassociating the three pullback comparisons, which is the pseudofunctor pentagon
  \texttt{pseudofunctor\_associativity}; the inner adjunction unit is slid past the
  outer pushforward by naturality of the unit (\(\mathrm{Adjunction.unit\_naturality}\)).

  The content is entirely Mathlib-side pseudofunctor coherence; the work is the
  bookkeeping of the \(\mathrm{eqToHom}\) coercions along the over-triangle
  \(p_1 \circ \bar g = p_2\). No step invokes the project-local sheafification/pullback
  square, so the laws are independent of the tensor--pullback substrate.
\end{proof}

\begin{lemma}
[Base-change unit (mate) identity for the push--pull head]
  \label{lem:push_pull_unit_mate}
  \lean{AlgebraicGeometry.pushPull_unit_mate}
  \uses{def:push_pull_map}
  For composable scheme morphisms \(f : A \to B\), \(p : B \to Z\) and a
  \(\mathcal{O}_Z\)-module \(N\), the adjunction unit \(\eta^p_N\) of
  \(p^* \dashv p_*\) at \(N\), followed by \(p_*\) of the unit \(\eta^f\) for
  \(f^* \dashv f_*\) at \(p^*N\), followed by the pushforward comparison
  \(p_* f_* \cong (f \circ p)_*\), equals the unit \(\eta^{f \circ p}_N\) for the
  composite followed by \((f \circ p)_*\) applied to the inverse pullback
  comparison \((f \circ p)^* \cong f^* p^*\):
  \[
    \eta^p_N \;\circ\; p_*(\eta^f_{p^*N}) \;\circ\;
      \bigl(p_* f_* \cong (f \circ p)_*\bigr)
    \;=\;
    \eta^{f \circ p}_N \;\circ\; (f \circ p)_*\bigl((f\circ p)^* \cong f^* p^*\bigr)^{-1}.
  \]
  This is the mate-calculus core that converts the single-morphism unit
  \(\eta^{f \circ p}\) into the iterated units \(\eta^p, \eta^f\); it is the
  reusable ingredient consumed by the functoriality (pentagon) law of the
  push--pull comparison map of Definition~\ref{def:push_pull_map}.
\end{lemma}

\begin{proof}
  Proved directly in Lean.
\end{proof}

\begin{lemma}
[Over-triangle transport cancellation for the push--pull tail]
  \label{lem:push_pull_transport_cancel}
  \lean{AlgebraicGeometry.pushPull_transport_cancel}
  \uses{def:push_pull_map}
  Let \(\bar g : Y_2 \to Y_1\), \(p_1 : Y_1 \to X\), \(p_2 : Y_2 \to X\) satisfy the
  over-triangle \(\bar g \circ p_1 = p_2\), and let \(\mathcal{F}\) be a
  \(\mathcal{O}_X\)-module. Then the pullback-comparison leg of the push--pull map,
  conjugated by the two \(\mathrm{eqToHom}\) coercions along the over-triangle,
  collapses to the transport-light form
  \[
    e_1 \;\circ\; p_{2*}\bigl((\bar g^* p_1^* \cong (\bar g p_1)^*)_{\mathcal{F}}\bigr)
      \;\circ\; e_2
    \;=\;
    (\bar g p_1)_*\bigl((\bar g^* p_1^* \cong (\bar g p_1)^*)_{\mathcal{F}}\bigr)
      \;\circ\; e_3,
  \]
  where \(e_1, e_2, e_3\) are the \(\mathrm{eqToHom}\) coercions induced by
  \(\bar g \circ p_1 = p_2\). It states the over-triangle cancellation in the
  push--pull comparison map of Definition~\ref{def:push_pull_map} with the
  over-triangle equality as a free hypothesis, so a single substitution turns the
  transports into identities; it is the reusable pre-coherence brick for the
  composition (pentagon) law.
\end{lemma}

\begin{proof}
  Proved directly in Lean.
\end{proof}

\begin{definition}\leanok
[Relative {\v C}ech complex from a nerve]
  \label{def:relative_cech_complex_of_nerve}
  \lean{AlgebraicGeometry.relativeCechComplexOfNerve}
  \uses{def:cech_nerve}
  Given \(f : X \to S\) and an augmented cosimplicial \(\mathcal{O}_X\)-module
  \(N\), the \emph{relative {\v C}ech cochain complex} in \(\operatorname{QCoh}(S)\)
  obtained by forgetting the augmentation of \(N\), pushing the resulting
  cosimplicial object forward degreewise along \(f\), and taking the
  alternating-coface-map cochain complex. This passage uses only the preadditivity
  of \(\operatorname{QCoh}(S)\) and \emph{no} pushforward/pullback coherence; it is
  the coherence-free plumbing (3) that turns the nerve of
  Definition~\ref{def:cech_nerve} into the {\v C}ech complex.
\end{definition}

\section{Affine acyclicity of the {\v C}ech complex}

\begin{lemma}

\leanok
[{\v C}ech acyclicity on affines]
  \label{lem:cech_acyclic_affine}
  \lean{AlgebraicGeometry.CechAcyclic.affine}
  \uses{def:cech_complex}
  % SOURCE: [Stacks Project], Cohomology of Schemes, Tag 02KG (affine vanishing);
  % Lemma \texttt{lemma-cech-cohomology-quasi-coherent-trivial} (standard-cover
  % {\v C}ech vanishing) and Lemma
  % \texttt{lemma-quasi-coherent-affine-cohomology-zero} (Serre vanishing on affines)
  % (read from references/stacks-coherent.tex, L44--52 and L145--155).
  % SOURCE QUOTE: "Let $X$ be a scheme.
  % Let $\mathcal{F}$ be a quasi-coherent $\mathcal{O}_X$-module.
  % Let $\mathcal{U} : U = \bigcup_{i = 1}^n D(f_i)$ be a standard
  % open covering of an affine open of $X$.
  % Then $\check{H}^p(\mathcal{U}, \mathcal{F}) = 0$ for
  % all $p > 0$."
  % SOURCE QUOTE: "Serre vanishing: Higher cohomology vanishes on affine schemes
  % for quasi-coherent modules.
  % Let $X$ be a scheme.
  % Let $\mathcal{F}$ be a quasi-coherent $\mathcal{O}_X$-module.
  % For any affine open $U \subset X$ we have
  % $H^p(U, \mathcal{F}) = 0$ for all $p > 0$."
  \textit{Source: Stacks Project, Cohomology of Schemes,
  \texttt{lemma-cech-cohomology-quasi-coherent-trivial} (standard-cover {\v C}ech
  vanishing) and Stacks, Tag 02KG
  (\texttt{lemma-quasi-coherent-affine-cohomology-zero}, Serre vanishing on
  affines).}
  Let \(U = \operatorname{Spec}(A)\) be an affine scheme, \(\mathcal{F}\) a
  quasi-coherent \(\mathcal{O}_U\)-module, and
  \(\mathcal{U} : U = \bigcup_{i=1}^n D(f_i)\) a standard open cover (so the
  \(f_i \in A\) generate the unit ideal). Then the {\v C}ech complex
  \(\check{\mathcal{C}}^\bullet(\mathcal{U}, \mathcal{F})\) has vanishing cohomology
  in all positive degrees:
  \[
    \check{H}^p(\mathcal{U}, \mathcal{F}) = 0 \qquad \text{for all } p > 0,
  \]
  and consequently \(H^p(U, \mathcal{F}) = 0\) for all \(p > 0\).
\end{lemma}

% SOURCE QUOTE PROOF: "Write $U = \Spec(A)$ for some ring $A$.
% In other words, $f_1, \ldots, f_n$ are elements of $A$
% which generate the unit ideal of $A$.
% Write $\mathcal{F}|_U = \widetilde{M}$ for some $A$-module $M$.
% ...
% It suffices to show that (\ref{equation-extended})
% is exact after localizing at a prime $\mathfrak p$, see
% Algebra, Lemma \ref{algebra-lemma-characterize-zero-local}.
% In fact we will show that the extended complex localized
% at $\mathfrak p$ is homotopic to zero.
% There exists an index $i$ such that $f_i \not \in \mathfrak p$.
% Choose and fix such an element $i_{\text{fix}}$. ...
% Let us define a homotopy
% $$
% h :
% \prod\nolimits_{i_0 \ldots i_{p + 1}}
% M_{f_{i_0} \ldots f_{i_{p + 1}}, \mathfrak p}
% \longrightarrow
% \prod\nolimits_{i_0 \ldots i_p}
% M_{f_{i_0} \ldots f_{i_p}, \mathfrak p}
% $$
% by the rule
% $$
% h(s)_{i_0 \ldots i_p} = s_{i_{\text{fix}} i_0 \ldots i_p}
% $$
% ... We compute ... $(dh + hd)(s)_{i_0 \ldots i_p} = s_{i_0 \ldots i_p}$ ...
% This proves the identity map is homotopic to zero as desired."
% SOURCE QUOTE PROOF: "We are going to apply
% Cohomology, Lemma \ref{cohomology-lemma-cech-vanish-basis}.
% As our basis $\mathcal{B}$ for the topology of $X$ we are going to use
% the affine opens of $X$. ... Note that the intersection of standard opens
% in an affine is another standard open. Hence property (1) holds. ...
% Finally, condition (3) of the lemma follows from
% Lemma \texttt{lemma-cech-cohomology-quasi-coherent-trivial}."
\begin{proof}
  \uses{def:cech_complex}
  Write \(U = \operatorname{Spec}(A)\) and \(\mathcal{F}|_U = \widetilde{M}\) for an
  \(A\)-module \(M\). By Definition~\ref{def:cech_complex} the {\v C}ech complex of
  the standard cover is the complex of localisations
  \[
    \prod_{i_0} M_{f_{i_0}} \;\to\;
    \prod_{i_0 i_1} M_{f_{i_0} f_{i_1}} \;\to\;
    \prod_{i_0 i_1 i_2} M_{f_{i_0} f_{i_1} f_{i_2}} \;\to\; \cdots,
  \]
  and the claim \(\check{H}^p = 0\) for \(p > 0\) is equivalent to exactness of the
  extended complex \(0 \to M \to \prod_{i_0} M_{f_{i_0}} \to \cdots\). Exactness may
  be checked after localising at an arbitrary prime \(\mathfrak{p} \subset A\).
  Choose an index \(i_{\mathrm{fix}}\) with
  \(f_{i_{\mathrm{fix}}} \notin \mathfrak{p}\); then
  \(M_{f_{i_{\mathrm{fix}}}, \mathfrak{p}} = M_{\mathfrak{p}}\), and the
  prescription \(h(s)_{i_0 \ldots i_p} = s_{i_{\mathrm{fix}} i_0 \ldots i_p}\)
  defines a contracting homotopy of the localised extended complex: a direct
  computation gives \((dh + hd)(s) = s\), so the identity is homotopic to zero and
  the localised complex is exact. As exactness holds after every localisation, the
  extended complex is exact and \(\check{H}^p(\mathcal{U}, \mathcal{F}) = 0\) for
  \(p > 0\).

  For the sheaf-cohomology statement, take the affine opens of \(U\) as a basis of
  the topology and the standard open covers as the admissible covers. Intersections
  of standard opens are again standard opens, standard covers are cofinal among
  covers of a basic open, and the {\v C}ech vanishing just proved supplies the
  acyclicity hypothesis; the {\v C}ech-to-cohomology comparison on a basis then
  yields \(H^p(U, \mathcal{F}) = 0\) for all \(p > 0\). \emph{(This is Serre
  vanishing for quasi-coherent sheaves on affine schemes.)}

  \emph{This proof depends on the following currently-absent Mathlib
  infrastructure: the explicit description of the standard-cover {\v C}ech complex
  as the complex of localisations \(\prod M_{f_{i_0}} \to \prod M_{f_{i_0}f_{i_1}}
  \to \cdots\) together with the prime-local, module-level contracting homotopy
  \(h(s)_{i_0 \ldots i_p} = s_{i_{\mathrm{fix}} i_0 \ldots i_p}\) witnessing its
  positive-degree exactness, neither of which is available for sheaves of modules
  on a scheme.}
\end{proof}

\section{The {\v C}ech complex computes $R^i f_*$}

\begin{lemma}

\leanok
[{\v C}ech complex computes the higher direct images]
  \label{lem:cech_computes_cohomology}
  \lean{AlgebraicGeometry.cech_computes_higherDirectImage}
  \uses{lem:cech_acyclic_affine, def:cech_complex}
  % SOURCE: [Stacks Project], Cohomology of Schemes, Tag 02KE; Lemma
  % \texttt{lemma-cech-cohomology-quasi-coherent} ({\v C}ech computes cohomology when
  % all intersections are affine) and Lemma
  % \texttt{lemma-quasi-coherence-higher-direct-images-application}
  % ($H^q(X,\mathcal{F}) = H^0(S, R^q f_*\mathcal{F})$ for affine $S$)
  % (read from references/stacks-coherent.tex, L245--256 and L843--854).
  % SOURCE QUOTE: "Let $X$ be a scheme.
  % Let $\mathcal{U} : X = \bigcup_{i \in I} U_i$ be an open covering such that
  % $U_{i_0 \ldots i_p}$ is affine open for all $p \ge 0$ and all
  % $i_0, \ldots, i_p \in I$.
  % In this case for any quasi-coherent sheaf $\mathcal{F}$ we have
  % $$
  % \check{H}^p(\mathcal{U}, \mathcal{F}) = H^p(X, \mathcal{F})
  % $$
  % as $\Gamma(X, \mathcal{O}_X)$-modules for all $p$."
  % SOURCE QUOTE: "Let $f : X \to S$ be a morphism of schemes.
  % Assume that $f$ is quasi-separated and quasi-compact.
  % Assume $S$ is affine.
  % For any quasi-coherent $\mathcal{O}_X$-module $\mathcal{F}$
  % we have
  % $$
  % H^q(X, \mathcal{F}) = H^0(S, R^qf_*\mathcal{F})
  % $$
  % for all $q \in \mathbf{Z}$."
  \textit{Source: Stacks Project, Cohomology of Schemes, Tag 02KE
  (\texttt{lemma-cech-cohomology-quasi-coherent}) and
  \texttt{lemma-quasi-coherence-higher-direct-images-application}.}
  Let \(f : X \to S\) be a separated quasi-compact morphism of schemes, let
  \(\mathcal{F}\) be a quasi-coherent \(\mathcal{O}_X\)-module, and let
  \(\mathfrak{U} : X = \bigcup_{i \in I} U_i\) be a finite affine open cover of
  \(X\) (so, by separatedness, every intersection \(U_{i_0 \ldots i_p}\) is
  affine). Then the cohomology sheaves of the relative {\v C}ech complex
  \(\check{\mathcal{C}}^\bullet(\mathfrak{U}, \mathcal{F})\) of
  Definition~\ref{def:cech_complex} compute the higher direct images: for every
  \(i \geq 0\) there is a canonical isomorphism of \(\mathcal{O}_S\)-modules
  \[
    H^i\bigl(\check{\mathcal{C}}^\bullet(\mathfrak{U}, \mathcal{F})\bigr)
      \;\cong\; R^i f_*\mathcal{F}.
  \]
  In particular, over an affine base \(S = \operatorname{Spec}(A)\), taking global
  sections gives \(H^i(X, \mathcal{F}) = \check{H}^i(\mathfrak{U}, \mathcal{F}) =
  H^0(S, R^i f_*\mathcal{F})\) as \(A\)-modules for all \(i\).

  Two points about the precise form of the comparison should be noted. First, the
  derived-functor higher direct image \(R^i f_*\mathcal{F}\) on the right-hand side
  is the one obtained from injective resolutions, which is available only when the
  category of \(\mathcal{O}_X\)-modules has enough injectives; accordingly this
  comparison is asserted under the hypothesis that injective resolutions exist in
  \(\operatorname{QCoh}(X)\), and it compares the (always-defined) {\v C}ech higher
  direct image of Definition~\ref{def:cech_higher_direct_image} against the
  derived-functor one only in that case. Second, we assert the comparison in the
  weak form of the \emph{existence} of an isomorphism, i.e.\ as
  \(\bigl(H^i(\check{\mathcal{C}}^\bullet(\mathfrak{U}, \mathcal{F})) \cong
  R^i f_*\mathcal{F}\bigr)\) without singling out a chosen natural isomorphism.
  This existence statement is all the downstream consumers (representability of the
  relative Picard functor, local triviality of the comparison) require: they use
  only that the unconditional {\v C}ech \(R^i f_*\) and the derived-functor
  \(R^i f_*\) agree as objects, never a specific identification, so committing to a
  canonical natural transformation would add coherence obligations that play no
  role in the applications.
\end{lemma}

% SOURCE QUOTE PROOF: "In view of
% Lemma \texttt{lemma-quasi-coherent-affine-cohomology-zero}
% this is a special case of
% Cohomology, Lemma
% \ref{cohomology-lemma-cech-spectral-sequence-application}."
% SOURCE QUOTE PROOF: "Consider the Leray spectral sequence
% $E_2^{p, q} = H^p(S, R^qf_*\mathcal{F})$
% converging to $H^{p + q}(X, \mathcal{F})$, see
% Cohomology, Lemma \ref{cohomology-lemma-Leray}.
% By Lemma \ref{lemma-quasi-coherence-higher-direct-images}
% we see that the sheaves $R^qf_*\mathcal{F}$ are quasi-coherent.
% By Lemma \texttt{lemma-quasi-coherent-affine-cohomology-zero}
% we see that $E_2^{p, q} = 0$ when $p > 0$.
% Hence the spectral sequence degenerates at $E_2$ and we win."
\begin{proof}
  \uses{lem:cech_acyclic_affine, def:cech_complex}
  The question is local on \(S\), and the formation of both sides commutes with
  restriction to opens of \(S\), so we may assume \(S = \operatorname{Spec}(A)\) is
  affine. Because \(f\) is separated and quasi-compact, each finite intersection
  \(U_{i_0 \ldots i_p}\) is affine, and \(f\) restricted to it is an affine
  morphism; by the affine acyclicity of Lemma~\ref{lem:cech_acyclic_affine} the
  higher direct images of \(\mathcal{F}\) over each such intersection vanish, so the
  terms of the relative {\v C}ech complex carry no hidden higher cohomology.

  Consider the {\v C}ech-to-cohomology spectral sequence of the cover
  \(\mathfrak{U}\), with \(E_2\)-page
  \(E_2^{p,q} = \check{H}^p(\mathfrak{U}, \underline{H}^q(\mathcal{F}))\)
  converging to \(H^{p+q}(X, \mathcal{F})\). Since all intersections are affine,
  Lemma~\ref{lem:cech_acyclic_affine} forces the presheaf cohomology
  \(\underline{H}^q(\mathcal{F})\) to vanish on the intersections for \(q > 0\), so
  the spectral sequence collapses to its row \(q = 0\): the {\v C}ech cohomology of
  \(\mathcal{F}\) for the cover \(\mathfrak{U}\) computes the sheaf cohomology,
  \(\check{H}^p(\mathfrak{U}, \mathcal{F}) = H^p(X, \mathcal{F})\) for all \(p\),
  and the same identification holds relatively, i.e.\ the cohomology sheaves of
  \(\check{\mathcal{C}}^\bullet(\mathfrak{U}, \mathcal{F})\) compute the cohomology
  sheaves of \(Rf_*\mathcal{F}\).

  It remains to identify these cohomology sheaves with \(R^i f_*\mathcal{F}\). Over
  the affine base \(S = \operatorname{Spec}(A)\), the Leray spectral sequence
  \(E_2^{p,q} = H^p(S, R^q f_*\mathcal{F}) \Rightarrow H^{p+q}(X, \mathcal{F})\)
  has \(R^q f_*\mathcal{F}\) quasi-coherent, whence \(H^p(S, R^q f_*\mathcal{F}) = 0\)
  for \(p > 0\) by Serre vanishing (Lemma~\ref{lem:cech_acyclic_affine}). The Leray
  sequence therefore degenerates at \(E_2\), giving
  \(H^i(X, \mathcal{F}) = H^0(S, R^i f_*\mathcal{F})\). Combining with the previous
  paragraph, \(H^i\bigl(\check{\mathcal{C}}^\bullet(\mathfrak{U}, \mathcal{F})\bigr)
  \cong R^i f_*\mathcal{F}\) as \(\mathcal{O}_S\)-modules for all \(i \geq 0\).

  \emph{This proof depends on the following currently-absent Mathlib
  infrastructure: the {\v C}ech-to-derived-functor spectral sequence and the Leray
  spectral sequence for the category of sheaves of modules on a scheme, neither of
  which is presently available there; the argument uses both to collapse onto the
  comparison isomorphism.}
\end{proof}

\section{The unconditional higher direct image}

\begin{definition}

\leanok
[Unconditional higher direct image via {\v C}ech]
  \label{def:cech_higher_direct_image}
  \lean{AlgebraicGeometry.cechHigherDirectImage}
  \uses{lem:cech_computes_cohomology}
  % SOURCE: [Stacks Project], Cohomology of Schemes, Lemma
  % \texttt{lemma-quasi-coherence-higher-direct-images-application}
  % (read from references/stacks-coherent.tex, L843--854); the unconditional
  % packaging via the {\v C}ech complex is an Archon-original construction grounded
  % in this lemma and Lemma \texttt{lemma-cech-cohomology-quasi-coherent}.
  % SOURCE QUOTE: "Let $f : X \to S$ be a morphism of schemes.
  % Assume that $f$ is quasi-separated and quasi-compact.
  % Assume $S$ is affine.
  % For any quasi-coherent $\mathcal{O}_X$-module $\mathcal{F}$
  % we have
  % $$
  % H^q(X, \mathcal{F}) = H^0(S, R^qf_*\mathcal{F})
  % $$
  % for all $q \in \mathbf{Z}$."
  \textit{Source: Stacks Project, Cohomology of Schemes,
  \texttt{lemma-quasi-coherence-higher-direct-images-application}; unconditional
  {\v C}ech packaging is Archon-original.}
  Let \(f : X \to S\) be a separated quasi-compact morphism and \(\mathcal{F}\) a
  quasi-coherent \(\mathcal{O}_X\)-module. Fix a finite affine open cover
  \(\mathfrak{U}\) of \(X\). The \emph{(unconditional) \(i\)-th higher direct
  image} is defined, for each \(i \geq 0\), as the \(i\)-th cohomology sheaf of the
  relative {\v C}ech complex,
  \[
    R^i f_*\mathcal{F} \;:=\;
      H^i\bigl(\check{\mathcal{C}}^\bullet(\mathfrak{U}, \mathcal{F})\bigr)
      \in \operatorname{QCoh}(S).
  \]
  This requires \emph{no} hypothesis that the category of \(\mathcal{O}_X\)-modules
  have enough injectives: the right-hand side is the cohomology of an explicit
  complex of quasi-coherent sheaves. By
  Lemma~\ref{lem:cech_computes_cohomology} the result is canonically isomorphic to
  the derived-functor higher direct image wherever the latter is defined, and in
  particular is independent of the chosen affine cover \(\mathfrak{U}\) up to
  canonical isomorphism: two affine covers admit a common affine refinement, and
  the induced refinement maps on {\v C}ech complexes are quasi-isomorphisms because
  both compute the same \(R^i f_*\mathcal{F}\) by
  Lemma~\ref{lem:cech_computes_cohomology}. For \(i = 0\) one recovers the ordinary
  pushforward \(R^0 f_*\mathcal{F} = f_*\mathcal{F}\).
\end{definition}

\section{Flat base change for the {\v C}ech higher direct images}

\begin{lemma}

\leanok
[Flat base change of $R^i f_*$ via {\v C}ech]
  \label{lem:cech_flat_base_change}
  \lean{AlgebraicGeometry.cech_flatBaseChange}
  \uses{def:cech_higher_direct_image, lem:cech_computes_cohomology}
  % SOURCE: [Stacks Project], Cohomology of Schemes, Tag 02KH (Lemma
  % \texttt{lemma-flat-base-change-cohomology}, "Flat base change")
  % (read from references/stacks-coherent.tex, L947--970).
  % SOURCE QUOTE: "Consider a cartesian diagram of schemes
  % $$
  % \xymatrix{
  % X' \ar[d]_{f'} \ar[r]_{g'} & X \ar[d]^f \\
  % S' \ar[r]^g & S
  % }
  % $$
  % Let $\mathcal{F}$ be a quasi-coherent $\mathcal{O}_X$-module
  % with pullback $\mathcal{F}' = (g')^*\mathcal{F}$.
  % Assume that $g$ is flat and that $f$ is quasi-compact and quasi-separated.
  % For any $i \geq 0$
  % \begin{enumerate}
  % \item the base change map of
  % Cohomology, Lemma \ref{cohomology-lemma-base-change-map-flat-case}
  % is an isomorphism
  % $$
  % g^*R^if_*\mathcal{F} \longrightarrow R^if'_*\mathcal{F}',
  % $$
  % \item if $S = \Spec(A)$ and $S' = \Spec(B)$, then
  % $H^i(X, \mathcal{F}) \otimes_A B = H^i(X', \mathcal{F}')$.
  % \end{enumerate}"
  \textit{Source: Stacks Project, Cohomology of Schemes, Tag 02KH
  (\texttt{lemma-flat-base-change-cohomology}, ``Flat base change'').}
  Consider the cartesian square
  \[
    \begin{array}{ccc}
      X' & \xrightarrow{\;g'\;} & X \\[2pt]
      {\scriptstyle f'}\big\downarrow & & \big\downarrow{\scriptstyle f} \\[2pt]
      S' & \xrightarrow{\;g\;} & S
    \end{array}
  \]
  with \(f : X \to S\) separated and quasi-compact, \(\mathcal{F}\) a
  quasi-coherent \(\mathcal{O}_X\)-module, and
  \(\mathcal{F}' = (g')^*\mathcal{F}\). Assume \(g\) is flat. Then for every
  \(i \geq 0\) the canonical base-change map between the unconditional {\v C}ech
  higher direct images of Definition~\ref{def:cech_higher_direct_image} is an
  isomorphism
  \[
    g^*\bigl(R^i f_*\mathcal{F}\bigr) \;\xrightarrow{\;\sim\;}\;
      R^i f'_*\bigl((g')^*\mathcal{F}\bigr) = R^i f'_*\mathcal{F}'.
  \]
  Equivalently, when \(S = \operatorname{Spec}(A)\) and
  \(S' = \operatorname{Spec}(B)\) with \(A \to B\) flat, the comparison map of
  \(B\)-modules \(H^i(X, \mathcal{F}) \otimes_A B \to H^i(X', \mathcal{F}')\) is an
  isomorphism.
\end{lemma}

% SOURCE QUOTE PROOF: "Having said this, part (1) follows from part (2). Namely,
% part (1) is local on $S'$ and hence we may assume $S$
% and $S'$ are affine. ...
% In case $X$ is separated, choose an affine open covering
% $\mathcal{U} : X = U_1 \cup \ldots \cup U_t$ and recall that
% $$
% \check{H}^p(\mathcal{U}, \mathcal{F}) = H^p(X, \mathcal{F}),
% $$
% see
% Lemma \texttt{lemma-cech-cohomology-quasi-coherent}.
% If $\mathcal{U}_B : X_B = (U_1)_B \cup \ldots \cup (U_t)_B$ we obtain
% by base change, then it is still the case that each $(U_i)_B$ is affine
% and that $X_B$ is separated. ...
% We have the following relation between the {\v C}ech complexes
% $$
% \check{\mathcal{C}}^\bullet(\mathcal{U}_B, \mathcal{F}_B) =
% \check{\mathcal{C}}^\bullet(\mathcal{U}, \mathcal{F}) \otimes_A B
% $$
% as follows from
% Lemma \ref{lemma-affine-base-change}.
% Since $A \to B$ is flat, the same thing remains true on taking cohomology."
\begin{proof}
  \uses{def:cech_higher_direct_image, lem:cech_computes_cohomology}
  The assertion that the base-change map is an isomorphism is local on \(S'\), so we
  may assume \(S = \operatorname{Spec}(A)\) and \(S' = \operatorname{Spec}(B)\) with
  \(A \to B\) flat. By Lemma~\ref{lem:cech_computes_cohomology} both
  \(R^i f_*\mathcal{F}\) and \(R^i f'_*\mathcal{F}'\) are quasi-coherent, associated
  to the \(A\)-module \(H^i(X, \mathcal{F})\) and the \(B\)-module
  \(H^i(X', \mathcal{F}')\) respectively; since pullback along \(g\) corresponds to
  \(- \otimes_A B\) on modules, it suffices to prove that the comparison map
  \(H^i(X, \mathcal{F}) \otimes_A B \to H^i(X', \mathcal{F}')\) is an isomorphism.

  Choose a finite affine open cover \(\mathfrak{U} : X = U_1 \cup \dots \cup U_t\);
  by separatedness all intersections \(U_{i_0 \ldots i_p}\) are affine. Base
  changing the cover yields an affine open cover \(\mathfrak{U}_B\) of the
  again-separated \(X' = X_B\), with each \((U_{i_0 \ldots i_p})_B\) affine. By
  Lemma~\ref{lem:cech_computes_cohomology} the two {\v C}ech complexes compute the
  respective cohomologies. The affine base change for the \(i = 0\) direct image
  identifies each term of the base-changed {\v C}ech complex with the original term
  tensored over \(A\) with \(B\), giving an isomorphism of cochain complexes
  \[
    \check{\mathcal{C}}^\bullet(\mathfrak{U}_B, \mathcal{F}_B)
      \;\cong\;
      \check{\mathcal{C}}^\bullet(\mathfrak{U}, \mathcal{F}) \otimes_A B.
  \]
  Because \(A \to B\) is flat, \(- \otimes_A B\) is exact and commutes with the
  formation of cohomology of a complex; applying \(H^i\) for any \(i \geq 0\) yields
  \[
    H^i(X, \mathcal{F}) \otimes_A B \;\xrightarrow{\;\sim\;}\;
      H^i(X', \mathcal{F}').
  \]
  Translating back to sheaves on \(S' = \operatorname{Spec}(B)\) — using that the
  unconditional higher direct images are the quasi-coherent sheaves associated to
  these modules — gives the asserted isomorphism
  \(g^*(R^i f_*\mathcal{F}) \xrightarrow{\sim} R^i f'_*\mathcal{F}'\) for all
  \(i \geq 0\). The flatness of \(g\) (hence of \(g'\), by base change) is exactly
  what makes \(- \otimes_A B\) exact, and is therefore the precise hypothesis under
  which the {\v C}ech computation of Definition~\ref{def:cech_higher_direct_image}
  is compatible with base change.

  \emph{This proof depends on the following currently-absent Mathlib
  infrastructure: the term-wise affine base change of the {\v C}ech complex
  identifying \(\check{\mathcal{C}}^\bullet(\mathfrak{U}_B, \mathcal{F}_B) \cong
  \check{\mathcal{C}}^\bullet(\mathfrak{U}, \mathcal{F}) \otimes_A B\), and the
  exactness of \(- \otimes_A B\) (its commutation with the formation of cohomology
  of a complex) for the relevant categories of sheaves of modules, neither of which
  is presently available there.}
\end{proof}
